\section{Related Work}

\paragraph{Single-cell aptamer profiling.} Existing single-cell aptamer studies address complementary goals in cell discrimination, biological characterization, and probe discovery \citep{xu2024aptamer, wu2024efficient}. Apt-seq jointly measured aptamer binding and transcriptomes in individual cells, demonstrating discrimination between human Ramos and mouse 3T3 cells from their binding profiles \citep{delley2018combined}. \citep{wu2025high} developed an integrative strategy termed Aptomics to profile cell-surface proteins, glycans, and mRNA, characterizing tumor heterogeneity in clinical samples and glycosylation changes associated with T-cell differentiation. Along a complementary direction, SPARK-seq integrates CRISPR-based surface-protein perturbations with single-cell mRNA and aptamer sequencing to identify aptamer–target relationships and characterize binding kinetics \citep{luo2026spark}. Building on these profiling capabilities and biological applications, APTBench focuses on the predictive use of a fixed APT panel: evaluating cell-identity prediction under a specified RNA-derived reference taxonomy in patients excluded from model development.

\paragraph{Cell annotation and hierarchical prediction.} Cell-type annotation methods assign molecular profiles to reference identities through supervised classification or reference-based matching, as exemplified by CellTypist \citep{dominguez2022cross} and SingleR \citep{aran2019reference}. At larger scales, scTab \citep{fischer2024sctab} studied cross-tissue annotation using donor-wise data splits, providing a precedent for evaluating predictions on unseen donors. In parallel, foundation models such as scGPT \citep{cui2024scgpt} and Geneformer \citep{theodoris2023transfer} learn representations from large collections of single-cell transcriptomes and adapt them to downstream tasks through fine-tuning. To accommodate related cell populations and differing annotation resolutions, CHETAH \citep{de2019chetah} and scHPL \citep{michielsen2021hierarchical} organize reference populations into classification hierarchies. Hierarchical cross-entropy (HCE) further incorporates ancestor--descendant relationships between cell labels into the training objective \citep{cultrera2026improving}. Drawing on these methodological foundations, APTBench evaluates classical classifiers and hierarchy-aware reference models with APT inputs and analyzes prediction errors across and within lineages. Separately, its true-lineage oracle supplies the reference lineage at test time to diagnose remaining subtype discrimination, with training-derived conditional priors receiving the same lineage information as controls. 